\chapter{Introduction to Natural Language Processing}

Natural Language Processing (NLP) is the field of study that focuses on the interaction between computers and human language. Its goal is to enable machines to understand, interpret, and generate human language in a way that is both meaningful and useful.

NLP is inherently interdisciplinary, drawing from linguistics, computer science, and artificial intelligence. It addresses the challenge of bridging the gap between the structured logic of machines and the complexity of human communication, which is rich in ambiguity, context, and cultural nuances.

The study of NLP is compelling because it allows us to automate and scale the understanding of vast amounts of textual data, from social media posts to scientific literature. By doing so, it unlocks new possibilities for extracting insights, improving accessibility, and creating more intuitive human-computer interfaces. NLP powers applications that we use daily, such as translation tools, virtual assistants, and sentiment analysis systems, making it a cornerstone of modern technology and a key driver of innovation in how we interact with information.